%LTeX: language=fr-FR
\documentclass[12pt]{article}
\usepackage[margin=0.8in]{geometry}
\usepackage[utf8]{inputenc}
\usepackage{graphicx}
\usepackage[export]{adjustbox}
\usepackage{amsmath}
\usepackage{amsfonts}
\usepackage{amssymb}
\usepackage[version=4]{mhchem}
\usepackage{stmaryrd}
\usepackage{bbold}
%\usepackage[fallback]{xeCJK}
%\usepackage{polyglossia}
\usepackage{fontspec}
\usepackage{censor}
\usepackage{currfile}
%\setCJKmainfont{Noto Serif CJK TC}
\usepackage{enumitem}
\usepackage{tikz}
\usetikzlibrary {arrows.meta, quotes} 
\usepackage{tkz-tab}
\usepackage{pgfplots}
\pgfplotsset{compat=1.18}
\usepackage{amsthm}
\usepackage{tcolorbox}
\usepackage{array}
\usepackage{colortbl}
\usepackage{tabularray}
\usepackage{tabularx}
\usepackage{makecell}
\usepackage{esvect}
%\usepackage{changepage}   % for the adjustwidth environment
\usepackage{enumitem}
\usepackage{tasks}
\usepackage{fancyhdr}
\usepackage[french]{babel}
\usepackage{lastpage}

\tcbuselibrary{theorems}

\definecolor{GrayCen}{HTML}{d2d2d2}
\definecolor{GrayBG}{HTML}{d9d9d9}
\definecolor{BlueDef}{HTML}{104e8b}
\definecolor{GreenProp}{HTML}{025540}
\definecolor{RedMet}{HTML}{cd5556}
\definecolor{BlackSav}{HTML}{000000}
\def\censorcolor{GrayCen}
\let\svcensorrule\censorrule
\renewcommand\censorrule[1]{%
\textcolor{\censorcolor}{\svcensorrule{#1}}}

\newtcbtheorem
[]% init options
{definition}% name
{Définition}% title
{%
  colback=GrayBG,
  colframe=BlueDef,
  fonttitle=\bfseries,
}% options
{def}% prefix

\newtcbtheorem
[]% init options
{proposition}% name
{Proposition}% title
{%
  colback=GrayBG,
  colframe=GreenProp,
  fonttitle=\bfseries,
}% options
{prop}% prefix

\newtcbtheorem
[]% init options
{method}% name
{Méthode}% title
{%
  theorem name,
  colback=GrayBG,
  colframe=RedMet,
  fonttitle=\bfseries,
}% options
{met}% prefix

\newtcbtheorem
[]% init options
{savoir}% name
{Savoir-faire du chapitre}% title
{%
  theorem name,
  colback=GrayBG,
  colframe=BlackSav,
  fonttitle=\bfseries,
}% options
{sav}% prefix

\newtheoremstyle{break}
{\topsep}{\topsep}%
{}{}%
{\bfseries}{}%
{\newline}{}%
\theoremstyle{break}
\newtheorem*{example}{Exemple~:}
\newtheorem*{examples}{Exemples~:}
\newtheorem*{remark}{Remarque~:}
\newtheorem*{remarks}{Remarques~:}
% Variante sans retour à la ligne après « Exemple. »
\newtheoremstyle{inline}
{\topsep}{\topsep}%
{}{}%
{\bfseries}{}%
{ }{}%
\theoremstyle{inline}
\newtheorem*{example*}{Exemple~:}

\newenvironment{demo}[1][\proofname]{%
  \begin{proof}[#1]$ $\newline
  }{%
  \end{proof}
}

%\setmainfont{CMU Serif}name

\title{
  \vspace{4em}
  \centering
  \tcbset{colframe=BlackSav,colback=GrayBG}
  \tcbox[tikznode]{
    Chapitre 5 \\ Fonctions du second degré
  }
}

\author{}
\date{}

\pagestyle{empty}
\graphicspath{ {../Images/} }
\input{\currfiledir ../def.tex}
\def \Document {Cours}


\setlength\parindent{0pt}
%\setlist[itemize]{topsep=0pt}
\setlist{nolistsep}
%\setlist[itemize]{noitemsep}
%\setlist[enumerate]{noitemsep}

\tikzset{cross/.style={path picture={
  \draw[black]
    (path picture bounding box.south east)--(path picture bounding box.north west)
    (path picture bounding box.south west)--(path picture bounding box.north east);
}}}

\newcolumntype{C}[1]{>{\centering\arraybackslash}m{#1}}
\newcolumntype{L}[1]{>{\raggedright\arraybackslash}m{#1}}
\newcolumntype{N}{@{}m{0pt}@{}}

\newcolumntype{Y}{>{\centering\arraybackslash}X}

\fancypagestyle{firststyle}
{
  \fancyhf{}
  \fancyhead[L]{\Structure}
  \fancyhead[R]{\Niveau}
  \fancyfoot[L]{Alexandre Fernandez}
  \fancyfoot[R]{{\thepage\  / \pageref{LastPage}}}
   %\renewcommand{\headrulewidth}{0pt} % removes horizontal header line
}


\begin{document}

\maketitle
\thispagestyle{firststyle}

\tableofcontents

% \newcommand{\multihide}[1]{\mbox{\vphantom{#1}\colorbox{gray!50}{\parbox{\textwidth}{{#1}}}}}
%\newcommand{\multihide}[1]{\mbox{\vphantom{#1}\colorbox{gray!50}{\parbox{\textwidth}{{\phantom{#1}}}}}}
% On crée une boîte temporaire pour stocker le texte
\newsavebox{\hidebox}

% -----------
% Macros des trous : version élève si \ELEVE est défini à la compilation
% (recette « prof + élève »), sinon version prof.
\ifdefined\ELEVE
% Version élève : contenu invisible
\newcommand{\hideM}[1]{\mbox{\vphantom{$#1$}\smash{\colorbox{gray!50}{\phantom{$\displaystyle #1$}}}} }
\newcommand{\hideMW}[1]{\mbox{\vphantom{$#1$}\smash{\colorbox{white}{\phantom{$\displaystyle #1$}}}} }
\newcommand{\hide}[1]{\mbox{\vphantom{#1}\smash{\colorbox{gray!50}{\phantom{#1}}}}}
\NewEnviron{multihide}{%
  \colorbox{gray!50}{
    \begin{lrbox}{\hidebox}%
        \begin{minipage}{\linewidth}%
          \BODY
        \end{minipage}%
    \end{lrbox}%
    % On affiche un espace invisible ayant les dimensions exactes de la boîte
    \vrule width 0pt height \ht\hidebox depth \dp\hidebox\hskip\wd\hidebox%
  }
}
% Corrigé long dans un cadre blanc : zone vide de même hauteur
% Argument optionnel : espace ajouté sous le corrigé, ex. \begin{multihideW}[2em]
\NewEnviron{multihideW}[1][0pt]{%
  \par\noindent\fbox{%
    \begin{lrbox}{\hidebox}%
      \begin{minipage}{\dimexpr\linewidth-2\fboxsep-2\fboxrule}%
        \BODY\par\vspace{#1}%
      \end{minipage}%
    \end{lrbox}%
    \vrule width 0pt height \ht\hidebox depth \dp\hidebox\hskip\wd\hidebox%
  }
}
\newcommand{\hidetk}[1]{}
\else
% Version prof : contenu visible sur fond gris
\newcommand{\hideM}[1]{\mbox{\vphantom{$#1$}\smash{\colorbox{gray!50}{$\displaystyle #1$}}} }
\newcommand{\hideMW}[1]{\mbox{\vphantom{$#1$}\smash{\colorbox{white}{$\displaystyle #1$}}} }
\newcommand{\hide}[1]{\mbox{\vphantom{#1}\smash{\colorbox{gray!50}{#1}}}}
\NewEnviron{multihide}{%
  \colorbox{gray!50}{
    \begin{minipage}{\linewidth}%
      \BODY
    \end{minipage}%
  }
}
% Corrigé long dans un cadre blanc (argument optionnel : espace ajouté)
\NewEnviron{multihideW}[1][0pt]{%
  \par\noindent\fbox{%
    \begin{minipage}{\dimexpr\linewidth-2\fboxsep-2\fboxrule}%
      \BODY\par\vspace{#1}%
    \end{minipage}%
  }
}
\newcommand{\hidetk}[1]{#1}
\fi
% -----------

\newcommand{\wvec}[1]{\overrightarrow{#1}}
\newcommand{\ssi}[0]{\emph{ssi}\ }
\newcommand{\setR}[0]{$\mathbb{R}$}

\newpage

\pagestyle{fancy}

\fancyhf{} % clear existing header/footer entries
\fancyhead[R]{ {\Niveau}}
\fancyhead[L]{ {\Chapitre}}
%\fancyhead[R]{\large \textbf{\textbf{\Document\ (B) -- 1h}}}
\fancyfoot[L]{ A. Fernandez }
\fancyfoot[R]{{\thepage\  / \pageref{LastPage}}}

% Repère quadrillé : \repere{xmin}{xmax}{ymin}{ymax}
\newcommand{\repere}[4]{%
  \draw[very thin, gray!40, step=1] (#1,#3) grid (#2,#4);
  \draw[->] (#1,0) -- ({#2+0.4},0) node[below right] {\small $x$};
  \draw[->] (0,#3) -- (0,{#4+0.4}) node[above left] {\small $y$};
  \foreach \x in {#1,...,#2} {\ifnum\x=0\else\draw (\x,0.08) -- (\x,-0.08) node[below] {\scriptsize $\x$};\fi}
  \foreach \y in {#3,...,#4} {\ifnum\y=0\else\draw (0.08,\y) -- (-0.08,\y) node[left] {\scriptsize $\y$};\fi}
  \node[below left] at (0,0) {\scriptsize $0$};
}


% Repère avec graduations choisies : \repereG{xmin}{xmax}{ymin}{ymax}{graduations x}{graduations y}
\newcommand{\repereG}[6]{%
  \draw[very thin, gray!40, step=1] (#1,#3) grid (#2,#4);
  \draw[->] (#1,0) -- ({#2+0.4},0) node[below right] {\small $x$};
  \draw[->] (0,#3) -- (0,{#4+0.4}) node[above left] {\small $y$};
  \foreach \x in {#5} {\draw (\x,0.1) -- (\x,-0.1) node[below] {\scriptsize $\x$};}
  \foreach \y in {#6} {\draw (0.08,\y) -- (-0.08,\y) node[left] {\scriptsize $\y$};}
  \node[below left] at (0,0) {\scriptsize $0$};
}

\section{Forme factorisée et racines}

\begin{example}
  Dans l'activité 1, la hauteur du dauphin est $f(t) = -0{,}3\,t\,(t - 6)$. Le dauphin est au niveau de l'eau quand $f(t) = 0$, c'est-à-dire pour $t = \hideM{0}$ et $t = \hideM{6}$.
\end{example}

\begin{definition}{Fonction du second degré}{}
  \begin{multihide}
    Une fonction $f$ définie sur $\mathbb{R}$ est une \textbf{fonction polynôme du second degré} si elle peut s'écrire $f(x) = ax^2 + bx + c$, avec $a \neq 0$ (\textbf{forme développée}).

    Certaines s'écrivent aussi $f(x) = a\,(x - x_1)(x - x_2)$ : c'est la \textbf{forme factorisée}.

    Un nombre $r$ est une \textbf{racine} de $f$ si $f(r) = 0$.
  \end{multihide}
\end{definition}

\begin{proposition}{}{}
  \begin{multihide}
    Les racines de $f(x) = a\,(x - x_1)(x - x_2)$ sont $x_1$ et $x_2$.
  \end{multihide}
\end{proposition}

\begin{remark}
  Un produit est nul si et seulement si l'un de ses facteurs est nul : $f(x) = 0 \iff x - x_1 = 0$ ou $x - x_2 = 0$.
\end{remark}

\begin{examples}
  Donner les racines de chaque fonction.

  \medskip
  \begin{minipage}[t]{0.31\linewidth}
    $f(x) = 3(x - 2)(x - 7)$
    \begin{multihideW}[0.5em]
      Racines : $2$ et $7$.
    \end{multihideW}
  \end{minipage}
  \hfill
  \begin{minipage}[t]{0.31\linewidth}
    $g(x) = -2x(x - 5)$
    \begin{multihideW}[0.5em]
      $g(x) = -2(x - 0)(x - 5)$ : racines $0$ et $5$.
    \end{multihideW}
  \end{minipage}
  \hfill
  \begin{minipage}[t]{0.31\linewidth}
    $h(x) = (x + 4)(x - 1)$
    \begin{multihideW}[0.5em]
      $x + 4 = x - (-4)$ : racines $-4$ et $1$.
    \end{multihideW}
  \end{minipage}
\end{examples}

\begin{method}{Vérifier qu'un nombre est racine}{}
  On calcule l'image de ce nombre : il est racine si et seulement si son image vaut $0$.
\end{method}

\begin{example}
  Soit $f(x) = x^2 - 5x + 6$. Le nombre $2$ est-il racine de $f$~? Et le nombre $1$~?
  \begin{multihideW}[0.5em]
    $f(2) = 2^2 - 5 \times 2 + 6 = 4 - 10 + 6 = 0$ : $2$ est une racine de $f$.

    $f(1) = 1^2 - 5 \times 1 + 6 = 2 \neq 0$ : $1$ n'est pas une racine de $f$.
  \end{multihideW}
\end{example}

\newpage
\begin{method}{Étudier le signe d'une forme factorisée}{}
  On dresse un tableau de signes avec \textbf{une ligne par facteur} (méthode du chapitre 3 pour le signe de $ax + b$), puis on applique la règle des signes d'un produit sur la dernière ligne.
\end{method}

\begin{example}
  Dresser le tableau de signes de $f(x) = 2(x - 1)(x - 5)$.

  \begin{center}
    \begin{tikzpicture}
      \tkzTabInit[lgt=2,espcl=2.2]{$x$/0.8,$2$/0.8,$x - 1$/0.8,$x - 5$/0.8,$f(x)$/0.8}{$-\infty$,$1$,$5$,$+\infty$}
      \hidetk{
        \tkzTabLine{,+,t,+,t,+,}
        \tkzTabLine{,-,z,+,t,+,}
        \tkzTabLine{,-,t,-,z,+,}
        \tkzTabLine{,+,z,-,z,+,}
      }
    \end{tikzpicture}
  \end{center}
\end{example}

\begin{proposition}{}{}
  \begin{multihide}
    Une fonction $f(x) = a\,(x - x_1)(x - x_2)$ est \textbf{du signe de $a$ à l'extérieur des racines}, et du signe contraire entre les racines.
  \end{multihide}
\end{proposition}

\newpage
\section{Allure, axe de symétrie, sommet, variations}

\begin{proposition}{}{}
  \begin{multihide}
    La courbe d'une fonction du second degré est une \textbf{parabole}.
    \begin{itemize}
      \item Si $a > 0$, la parabole est tournée vers le haut.
      \item Si $a < 0$, la parabole est tournée vers le bas.
    \end{itemize}
    Elle admet un \textbf{axe de symétrie} vertical. Le point de la parabole situé sur cet axe est le \textbf{sommet}.
  \end{multihide}
\end{proposition}

\begin{center}
  \begin{tabular}{C{0.4\linewidth}C{0.4\linewidth}}
    $a > 0$ & $a < 0$ \\[0.3em]
    \begin{tikzpicture}[scale=0.5]
      \draw[->] (-3,0) -- (3.3,0); \draw[->] (0,-1.2) -- (0,3.5);
      \hidetk{\draw[blue, thick, domain=-1.8:1.8, samples=50] plot (\x,{\x*\x-0.5});
      \draw[dashed] (0,-0.5) -- (0,3.2); \fill (0,-0.5) circle (3pt) node[below right] {\scriptsize sommet};}
    \end{tikzpicture}
    &
    \begin{tikzpicture}[scale=0.5]
      \draw[->] (-3,0) -- (3.3,0); \draw[->] (0,-1.2) -- (0,3.5);
      \hidetk{\draw[blue, thick, domain=-1.8:1.8, samples=50] plot (\x,{-\x*\x+2.8});
      \draw[dashed] (0,-1) -- (0,2.8); \fill (0,2.8) circle (3pt) node[above right] {\scriptsize sommet};}
    \end{tikzpicture}
  \end{tabular}
\end{center}

\begin{proposition}{Axe de symétrie d'une forme factorisée}{}
  \begin{multihide}
    La parabole de $f(x) = a\,(x - x_1)(x - x_2)$ admet pour axe de symétrie la droite d'équation $x = \dfrac{x_1 + x_2}{2}$ : elle passe \textbf{au milieu des deux racines}.

    Le sommet $S$ a pour abscisse $x_S = \dfrac{x_1 + x_2}{2}$ et pour ordonnée $f(x_S)$.
  \end{multihide}
\end{proposition}

\begin{proposition}{Variations}{}
  \begin{multihide}
    \begin{minipage}{0.48\linewidth}
      \centering
      Si $a > 0$ : minimum $f(x_S)$ en $x_S$.

      \smallskip
      \begin{tikzpicture}[scale=0.8]
        \tkzTabInit[lgt=1.2,espcl=1.6]{$x$/0.8,$f$/1.4}{$-\infty$,$x_S$,$+\infty$}
        \tkzTabVar{+/ ,-/$f(x_S)$,+/ }
      \end{tikzpicture}
    \end{minipage}
    \hfill
    \begin{minipage}{0.48\linewidth}
      \centering
      Si $a < 0$ : maximum $f(x_S)$ en $x_S$.

      \smallskip
      \begin{tikzpicture}[scale=0.8]
        \tkzTabInit[lgt=1.2,espcl=1.6]{$x$/0.8,$f$/1.4}{$-\infty$,$x_S$,$+\infty$}
        \tkzTabVar{-/ ,+/$f(x_S)$,-/ }
      \end{tikzpicture}
    \end{minipage}
  \end{multihide}
\end{proposition}

\begin{example}
  \begin{minipage}[t]{0.56\linewidth}
    \vspace{0pt}
    Étude de $f(x) = 2(x - 1)(x - 5)$.
    \begin{multihideW}[0.5em]
      Racines : $1$ et $5$.

      Axe de symétrie : $x = \dfrac{1 + 5}{2} = 3$.

      Sommet : $f(3) = 2 \times (3 - 1) \times (3 - 5) = 2 \times 2 \times (-2) = -8$, donc $S(3~;~-8)$.

      $a = 2 > 0$ : la parabole est tournée vers le haut, $f$ admet un minimum $-8$ en $3$.

      \smallskip
      \begin{tikzpicture}[scale=0.8]
        \tkzTabInit[lgt=1.2,espcl=1.8]{$x$/0.8,$f$/1.4}{$-\infty$,$3$,$+\infty$}
        \tkzTabVar{+/ ,-/$-8$,+/ }
      \end{tikzpicture}
    \end{multihideW}
  \end{minipage}
  \hfill
  \begin{minipage}[t]{0.4\linewidth}
    \vspace{0pt}
    \centering
    \begin{tikzpicture}[x=0.7cm, y=0.22cm]
      \repereG{-1}{7}{-9}{11}{2,4,6}{-8,-4,4,8}
      \hidetk{
        \draw[blue, thick, domain=-0.2:6.2, samples=60] plot (\x,{2*(\x-1)*(\x-5)});
        \draw[dashed, red] (3,-9) -- (3,11);
        \fill (3,-8) circle (2.5pt) node[below right] {\scriptsize $S$};
        \fill (1,0) circle (2.5pt) (5,0) circle (2.5pt);
      }
    \end{tikzpicture}
  \end{minipage}
\end{example}

\newpage
\section{Forme développée $ax^2 + bx + c$}

\begin{method}{Passer de la forme factorisée à la forme développée}{}
  On développe le produit (double distributivité), puis on réduit.
\end{method}

\begin{example}
  Développer $f(x) = 2(x - 1)(x - 5)$.
  \begin{multihideW}[0.5em]
    $f(x) = 2(x^2 - 5x - x + 5) = 2(x^2 - 6x + 5) = 2x^2 - 12x + 10$.
  \end{multihideW}
\end{example}

\begin{remark}
  Pour $f(x) = ax^2 + bx + c$, on a $f(0) = c$ : la parabole coupe l'axe des ordonnées au point $(0~;~c)$.
\end{remark}

\begin{method}{Trouver l'axe de symétrie d'une forme développée}{}
  \begin{enumerate}
    \item On résout l'équation $f(x) = c$ : elle se ramène à $ax^2 + bx = 0$, que l'on factorise par $x$.
    \item Ses deux solutions sont symétriques par rapport à l'axe : l'axe passe \textbf{au milieu} de ces deux solutions.
    \item On calcule l'ordonnée du sommet avec l'image de son abscisse.
  \end{enumerate}
  Aucune formule n'est à connaître.
\end{method}

\begin{example}
  \begin{minipage}[t]{0.62\linewidth}
    \vspace{0pt}
    Étude de $g(x) = -x^2 + 4x + 5$.
    \begin{multihideW}[0.5em]
      $g(0) = 5$. On résout $g(x) = 5$ :

      $-x^2 + 4x + 5 = 5 \iff -x^2 + 4x = 0 \iff x(-x + 4) = 0$

      $\iff x = 0$ ou $x = 4$.

      Axe de symétrie : $x = \dfrac{0 + 4}{2} = 2$.

      Sommet : $g(2) = -2^2 + 4 \times 2 + 5 = -4 + 8 + 5 = 9$, donc $S(2~;~9)$.

      $a = -1 < 0$ : la parabole est tournée vers le bas, $g$ admet un maximum $9$ en $2$.

      \smallskip
      \begin{tikzpicture}[scale=0.8]
        \tkzTabInit[lgt=1.2,espcl=1.8]{$x$/0.8,$g$/1.4}{$-\infty$,$2$,$+\infty$}
        \tkzTabVar{-/ ,+/$9$,-/ }
      \end{tikzpicture}
    \end{multihideW}
  \end{minipage}
  \hfill
  \begin{minipage}[t]{0.35\linewidth}
    \vspace{0pt}
    \centering
    \begin{tikzpicture}[x=0.6cm, y=0.25cm]
      \repereG{-2}{6}{-6}{10}{-1,2,4}{-4,4,8}
      \hidetk{
        \draw[blue, thick, domain=-1.5:5.5, samples=60] plot (\x,{-\x*\x+4*\x+5});
        \draw[dashed, red] (2,-6) -- (2,10);
        \fill (2,9) circle (2.5pt) node[above right] {\scriptsize $S$};
        \fill (0,5) circle (2.5pt) (4,5) circle (2.5pt);
        \draw[dotted] (0,5) -- (4,5);
      }
    \end{tikzpicture}
  \end{minipage}
\end{example}

\newpage
\section{Les fonctions $x \mapsto ax^2$ et $x \mapsto ax^2 + c$}

\begin{proposition}{}{}
  \begin{multihide}
    Pour $f(x) = ax^2 + c$ (cas où $b = 0$) : l'axe de symétrie est \textbf{l'axe des ordonnées} et le sommet est le point $S(0~;~c)$. Pour $f(x) = ax^2$, le sommet est l'origine du repère.
  \end{multihide}
\end{proposition}

\begin{method}{Associer une parabole à son expression}{}
  \begin{enumerate}
    \item Le signe de $a$ : parabole tournée vers le haut ($a > 0$) ou vers le bas ($a < 0$).
    \item L'image de $0$ : point d'intersection avec l'axe des ordonnées.
    \item Les racines : points d'intersection avec l'axe des abscisses.
  \end{enumerate}
\end{method}

\begin{example}
  Associer chaque parabole à son expression : $f_1(x) = x^2 - 4$, \quad $f_2(x) = -0{,}5x^2 + 2$, \quad $f_3(x) = (x - 1)(x - 3)$, \quad $f_4(x) = -(x + 1)(x - 3)$.

  \medskip
  \noindent
  \begin{tabular}{@{}C{0.23\linewidth}C{0.23\linewidth}C{0.23\linewidth}C{0.23\linewidth}@{}}
    \begin{tikzpicture}[scale=0.42]
      \repereG{-3}{4}{-4}{4}{-2,2}{-4,2}
      \draw[blue, thick, domain=-1.6:3.6, samples=50] plot (\x,{-(\x+1)*(\x-3)});
    \end{tikzpicture}
    &
    \begin{tikzpicture}[scale=0.42]
      \repereG{-3}{4}{-4}{4}{-2,2}{-4,2}
      \draw[blue, thick, domain=-2.45:2.45, samples=50] plot (\x,{\x*\x-4});
    \end{tikzpicture}
    &
    \begin{tikzpicture}[scale=0.42]
      \repereG{-3}{4}{-4}{4}{-2,2}{-4,2}
      \draw[blue, thick, domain=-0.3:4, samples=50] plot (\x,{(\x-1)*(\x-3)});
    \end{tikzpicture}
    &
    \begin{tikzpicture}[scale=0.42]
      \repereG{-3}{4}{-4}{4}{-2,2}{-4,2}
      \draw[blue, thick, domain=-3:3, samples=50] plot (\x,{-0.5*\x*\x+2});
    \end{tikzpicture}
    \\
    Courbe A : \hideM{f_4} & Courbe B : \hideM{f_1} & Courbe C : \hideM{f_3} & Courbe D : \hideM{f_2}
  \end{tabular}
\end{example}

\section{Factoriser connaissant une racine}

\begin{method}{Factoriser dans des cas simples}{}
  \begin{itemize}
    \item \textbf{Sans terme constant} ($c = 0$) : on factorise par $x$. Par exemple, $x^2 - 4x = x(x - 4)$.
    \item \textbf{Une racine $r$ est connue} : on écrit $f(x) = a\,(x - r)(x - k)$, où $a$ est le coefficient de $x^2$. On trouve $k$ en comparant les termes constants, puis on vérifie en développant.
  \end{itemize}
\end{method}

\begin{examples}
  \begin{enumerate}
    \item Factoriser $f(x) = x^2 - 5x + 6$, sachant que $2$ est une racine.
    \begin{multihideW}[0.5em]
      $f(x) = (x - 2)(x - k)$. Le terme constant de $(x - 2)(x - k)$ est $(-2) \times (-k) = 2k$, donc $2k = 6$ et $k = 3$.

      $f(x) = (x - 2)(x - 3)$. Vérification : $(x - 2)(x - 3) = x^2 - 3x - 2x + 6 = x^2 - 5x + 6$.
    \end{multihideW}

    \item Vérifier que $3$ est une racine de $h(x) = 2x^2 - 2x - 12$, puis factoriser $h(x)$.
    \begin{multihideW}[0.5em]
      $h(3) = 2 \times 9 - 2 \times 3 - 12 = 18 - 6 - 12 = 0$ : $3$ est une racine.

      $h(x) = 2(x - 3)(x - k)$. Terme constant : $2 \times (-3) \times (-k) = 6k$, donc $6k = -12$ et $k = -2$.

      $h(x) = 2(x - 3)(x + 2)$. Vérification : $2(x^2 + 2x - 3x - 6) = 2(x^2 - x - 6) = 2x^2 - 2x - 12$.
    \end{multihideW}
  \end{enumerate}
\end{examples}

\section{Résolution graphique et balayage}

\begin{method}{Résoudre graphiquement $f(x) = k$}{}
  On trace la droite horizontale d'équation $y = k$ : les solutions sont les \textbf{abscisses} des points d'intersection avec la courbe (chapitre 1).
\end{method}

\begin{method}{Encadrer une solution par balayage}{}
  Quand l'équation n'a pas de solution évidente, on calcule des images avec la table de la calculatrice : d'abord avec un pas de $1$, puis avec un pas de $0{,}1$ autour de la solution. Un solveur donne aussi une valeur approchée.
\end{method}

\begin{example}
  Le dauphin de l'activité 1 a pour hauteur $f(t) = -0{,}3t^2 + 1{,}8t$. À quels instants est-il à $2$~m de hauteur~?

  \renewcommand{\arraystretch}{1.3}
  \begin{center}
  \begin{tabular}{|c|*{4}{C{1.4cm}|}}
    \hline
    $t$ & $1$ & $2$ & $1{,}4$ & $1{,}5$ \\
    \hline
    $f(t)$ & $1{,}5$ & $2{,}4$ & $1{,}932$ & $2{,}025$ \\
    \hline
  \end{tabular}
  \end{center}
  \begin{multihideW}[0.5em]
    $f(1) < 2 < f(2)$, puis $f(1{,}4) < 2 < f(1{,}5)$ : une solution est comprise entre $1{,}4$ et $1{,}5$.

    Par symétrie par rapport à l'axe $t = 3$, l'autre solution est comprise entre $4{,}5$ et $4{,}6$.
  \end{multihideW}
\end{example}

\vfill

\begin{savoir}{}{}
\begin{itemize}
  \item Donner les racines d'une forme factorisée ; vérifier qu'un nombre est racine,
  \item Dresser le tableau de signes d'une forme factorisée,
  \item Déterminer l'axe de symétrie, le sommet et les variations (forme factorisée ou développée),
  \item Développer une forme factorisée ; trouver l'axe en résolvant $f(x) = c$,
  \item Associer une parabole à son expression,
  \item Factoriser dans des cas simples, connaissant une racine,
  \item Résoudre graphiquement $f(x) = k$ ; encadrer une solution par balayage.
\end{itemize}
\end{savoir}

\end{document}
